%%%
%
% $Autor: Siddhanth Sharma $
% $Datum: 2025-11-25 $
% $Dateiname: 
% $Version: 4620 $
%
% !TeX spellcheck = GB
% !TeX program = pdflatex
% !TeX encoding = utf8
%
%%%
\newpage
\section{\textcolor{red}{OpenCV}}

\index{OpenCV}

% --------------------------------------------------------
\subsection{Introduction} 
\index{OpenCV!Introduction}

OpenCV (Open Source Computer Vision Library) is an open-source library for real-time computer vision and image processing, originally developed by Intel and now maintained by the non-profit organisation OpenCV.org \cite{Bradski:2000,opencv_wiki}. It provides a large collection of optimised algorithms for working with images and videos, including filtering, geometric transformations, feature extraction, object detection, and tracking.

In this project, OpenCV is used as the core image-processing backbone for the licence plate detection pipeline on the Jetson Nano B01. It is responsible for:
\begin{itemize}
	\item capturing frames from the camera module via \PYTHON{cv2.VideoCapture},
	\item pre-processing images (resizing, colour conversion, normalisation) before inference,
	\item drawing bounding boxes, labels, and confidence scores on detected licence plates, and
	\item optionally displaying or saving annotated frames for debugging and evaluation.
\end{itemize}
The Python bindings (OpenCV-Python) allow seamless integration of these operations into the YOLOv8-based detection workflow \cite{opencv_py_intro}.


% --------------------------------------------------------
\subsection{Description}
\index{OpenCV!Description}

OpenCV is structured into several modules. Relevant modules include:
\begin{itemize}
	\item \textbf{core}: Provides functions for data structures such as matrices, random number generators.
	\item \textbf{imgproc}: Aids in image processing tasks like blurring, edge detection, morphological operations, and geometric transformations.
	\item \textbf{highgui}: GUI functions that enable displaying images and videos in windows.
	\item \textbf{videoio}: Enables reading and writing operations from video streams and cameras.
	\item \textbf{dnn}: Provides functionality for model inference along with loading and executing pre-trained models.
\end{itemize}

For the licence plate detection system, the typical processing chain is:
\begin{enumerate}
	\item Frame acquisition from the camera interface connected to the Jetson Nano using \PYTHON{cv2.VideoCapture}.
	\item Frame conversion from BGR (OpenCV default) to RGB or grayscale using \PYTHON{cv2.cvtColor}.
	\item Frame resizing as required by the YOLOv8 model (e.g.\ $640 \times 640$ pixels).
	\item Post-process model predictions. Rectangular and textual overlay on detected licence plates using \PYTHON{cv2.rectangle} and \PYTHON{cv2.putText}.
	\item Frames are saved and displayed for further processing.
\end{enumerate}

OpenCV is highly optimised and takes advantage of hardware acceleration (e.g.\ NEON, CUDA, and other backends) on embedded platforms. This makes it suitable for real-time processing on devices such as the Jetson Nano \cite{opencv_lfs,opencv_on_wheels}.


% --------------------------------------------------------
\subsection{Installation}
\index{OpenCV!Installation}

For this project, OpenCV is used via its Python bindings (\PYTHON{cv2}) in the JetPack~5.x environment on the Jetson Nano B01. There are two common installation strategies:

\begin{itemize}
	\item \textbf{Pre-built Python wheels} from PyPI, e.g.\ \texttt{opencv-python}, which provide CPU-only OpenCV binaries and are easy to install with \PYTHON{pip} \cite{opencv_pypi}.
	\item \textbf{Distribution or JetPack packages}, e.g.\ \texttt{python3-opencv} installed via \texttt{apt}, which are often optimised for the underlying hardware.
\end{itemize}

In the context of this project, OpenCV is installed using \PYTHON{pip} within the Jetson Nano’s virtual environment:

\begin{tcolorbox}[colback=white,colframe=black,coltext= ShellColor, sharp corners, boxrule=0.8pt]
	\begin{small}
		\begin{verbatim}
			
			pip install opencv-python	\end{verbatim}
	\end{small}
\end{tcolorbox}


%\SHELL{pip install opencv-python}

The command is responsible for installation of the \PYTHON{cv2} module, imported in Python scripts. For more advanced configurations (e.g.\ custom builds with CUDA support), the official OpenCV installation and configuration tutorials provide detailed CMake-based build instructions \cite{opencv_config_ref,opencv_py_tutorial_root}.


% --------------------------------------------------------
\subsection{Example -- Description}
\index{OpenCV!Example Description}

The example script \FILE{OpenCVExample.py} describes a minimal OpenCV workflow. The following points are verified:
\begin{itemize}
	\item The camera interface is initialized correctly and accessible,
	\item OpenCV is able to capture and display frames, and
	\item Image-processing operations function as intended.
\end{itemize}

In terms of concept, the following steps are performed:
\begin{enumerate}
	\item The default camera device loads using \PYTHON{cv2.VideoCapture(0)}.
	\item Frames are then read in a loop and valid frames returned. 
	\item Greyscale conversion using \PYTHON{cv2.cvtColor(frame, cv2.COLOR\_BGR2GRAY)} for preprocessing simulation prior to model inference.
	\item Static rectangular region of interest (ROI) generation where a license plate would be located, using \PYTHON{cv2.rectangle}.
	\item Processed frames are displayed in a window \PYTHON{cv2.imshow}.
	\item Loop termination on specific key presses (e.g.\ \texttt{q}) and resource release.
\end{enumerate}

Although this example does not run the YOLOv8 model itself, it mirrors the essential I/O and pre-processing steps used in the full licence plate detection pipeline.


% --------------------------------------------------------
\subsection{Manual Example}
\index{OpenCV!Example Manual}

The following set of instructions lay out running the OpenCV example on the Jetson Nano:

\begin{itemize}
	\item \textbf{Step 1:} Ensure the Jetson Nano is active and the camera is connected to the correct USB port.
	\item \textbf{Step 2:} Activate the Python virtual environment (if used) containing project dependencies, including OpenCV.
	\item \textbf{Step 3:} Navigate to the project directory where \FILE{OpenCVExample.py} is stored, for example:

	\begin{tcolorbox}[colback=white,colframe=black,coltext= ShellColor, sharp corners, boxrule=0.8pt]
		\begin{small}
			\begin{verbatim}
				
				cd ~/projects/license-plate-detection/opencv_tests	\end{verbatim}
		\end{small}
	\end{tcolorbox}

%	\begin{verbatim}
%		cd ~/projects/license-plate-detection/opencv_tests
%	\end{verbatim}


	\item \textbf{Step 4:} Run the example script:
		\begin{tcolorbox}[colback=white,colframe=black,coltext= ShellColor, sharp corners, boxrule=0.8pt]
		\begin{small}
			\begin{verbatim}
				
				python OpenCVExample.py	\end{verbatim}
		\end{small}
	\end{tcolorbox}

%	\begin{verbatim}
%		python OpenCVExample.py
%	\end{verbatim}


	\item \textbf{Step 5:} A window opens displaying the camera feed consisting of a grayscale conversion. 
	\item \textbf{Step 6:} Press the configured exit key (e.g.\ \texttt{q}) to close the window, release the camera, and terminate the program.
\end{itemize}

This procedure is vital for checking camera alignment and image quality prior to running full YOLOv8-based detection script.


% --------------------------------------------------------
\subsection{Example -- Version}
\index{OpenCV!Example Version}

In order to enact reproducibility and enable documentation of the exact OpenCV build implemented on the Jetson Nano, the following Python snippet may be used. The script is responsible for relaying the installed OpenCV package version.

\lstinputlisting[
language=Python,
caption={Checking the installed OpenCV version},
label={lst:opencv_version_check},
basicstyle=\ttfamily\small,
keywordstyle=\color{blue},
commentstyle=\color{gray},
stringstyle=\color{green!40!black},
breaklines=true,
showstringspaces=false,
numbers=left,
numberstyle=\tiny\color{gray},
frame=single,
rulecolor=\color{black!30}
]{../../Code/OpenCV/opencvversion.py}

Recording the version number in the project documentation helps to avoid compatibility issues when upgrading JetPack or when reinstalling packages on another device. For example, the online documentation pages used in this project correspond to OpenCV version 4.12.0, which provides the \emph{OpenCV-Python Tutorials} and related examples for image and video processing \cite{opencv_py_tutorial_root}.


% --------------------------------------------------------
\subsection{Example -- Files}
\index{OpenCV!Example Files}

The central example file for this section is:

\FILE{opencvexample.py}

This script is kept intentionally short and focused on:
\begin{itemize}
	\item initialising the camera,
	\item basic frame pre-processing, and
	\item visual verification of the image stream.
\end{itemize}
In the final system, similar code fragments are integrated into the main application script that loads the YOLOv8 model, performs inference on each frame, and writes detection results to the SQLite database.

\lstinputlisting[
language=Python,
caption={OpenCV Camera Test Example for Jetson Nano},
label={lst:opencv_example},
basicstyle=\ttfamily\small,
keywordstyle=\color{blue},
commentstyle=\color{gray},
stringstyle=\color{green!40!black},
breaklines=true,
showstringspaces=false,
numbers=left,
numberstyle=\tiny\color{gray},
frame=single,
rulecolor=\color{black!30}
]{../../Code/OpenCV/opencvexample.py}

\subsection{Further Reading}
\index{YOLO!Further Reading}

\begin{itemize}
	\item Python bindings (OpenCV-Python) and their integration into computer vision workflows \cite{opencv_py_intro}.
	\item Modular OpenCV architecture and optimised image and video processing pipelines \cite{opencv_lfs}.
	\item Hardware-accelerated OpenCV usage on embedded and edge platforms such as Jetson devices \cite{opencv_on_wheels}.
	\item Installation methods and distribution of OpenCV Python wheels via PyPI \cite{opencv_pypi}.
	\item Advanced OpenCV configuration, custom builds, and CMake-based compilation with optional CUDA support \cite{opencv_config_ref}.
\end{itemize}